%%%%%%%%%%%%%%%%%%%%%%%%%%%%%%%%%
% P: INTRO
%%%%%%%%%%%%%%%%%%%%%%%%%%%%%%%%%
%\ifthenelse{\equal{\deutsch}{false}}
%{\part{Introduction-Part}}
%{\part{Einleitung}}

%%%%%%%%%%%%%%%%%%%%%%%%%%%%%%%%%
% C: Requirements
%%%%%%%%%%%%%%%%%%%%%%%%%%%%%%%%%
\chapter{Requirements}
\input{\TEXTa 01_requirements_tab01}

%%%%%%%%%%%%%%%%%%%%%%%%%%%%%%%%%
% C: Document Overview
%%%%%%%%%%%%%%%%%%%%%%%%%%%%%%%%%
\chapter{Document Overview}
\input{\TEXTa 01_docover_text01}
\section{Validity of this document}
\input{\TEXTa 01_docover_text02}
\section{Target Specification Status}
\input{\TEXTa 01_docover_text03}
\section{Major Changes since last Revision}
\input{\TEXTa 01_docover_history_tab01}
%\section{Table of Contents}
\tableofcontents
%\section{List of Figures}
\listoffigures
%\section{List of Tables}
\listoftables
\section{Glossary}
\input{\TEXTa 01_docover_tab01}
\section{Naming Conventions}
\input{\TEXTa 01_docover_tab02}
\section{Document List}
\input{\TEXTa 01_docover_history_tab02}
\input{\TEXTa 01_docover_tab03}

%%%%%%%%%%%%%%%%%%%%%%%%%%%%%%%%%
% C: Product Overview
%%%%%%%%%%%%%%%%%%%%%%%%%%%%%%%%%
\chapter{Product Overview}
\section{Top Level View}
\input{\TEXTprod 02_product_history_tab01}
\subsection{Introduction}
\input{\TEXTprod 02_product_text01}
\input{\TEXTprod 02_product_bild01}
\subsection{Key Features}
\input{\TEXTprod 02_product_text03}
\subsection{Functional Block Diagram}
\input{\TEXTprod 02_product_text04}
\input{\TEXTprod 02_product_bild02}
\input{\TEXTprod 02_product_bild03}
\subsection{Package}
\input{\TEXTprod 02_product_text05}
\section{System View}
\input{\TEXTprod 02_product_history_tab02}
\input{\TEXTprod 02_product_text06}
\section{Architecture Concepts Overview}
\input{\TEXTprod 02_product_history_tab03}
\subsection{Technology}
\input{\TEXTprod 02_product_text07}
\subsection{System Memory Concept}
\input{\TEXTprod 02_product_text08}
\subsection{Software Architecture}
\input{\TEXTprod 02_product_text09}
\subsection{Interprocessor Communication Concept}
\input{\TEXTprod 02_product_text10}
\subsection{Bus Concept}
\input{\TEXTprod 02_product_text11}
\subsection{Clock Concept}
\input{\TEXTprod 02_product_text12}
\subsection{System Control Concept}
\input{\TEXTprod 02_product_text13}
\subsection{Interrupt Concept}
\input{\TEXTprod 02_product_text14}
\subsection{Debug Concept}
\input{\TEXTprod 02_product_text15}
\subsection{Power Management}
\input{\TEXTprod 02_product_text16}
\subsection{Protection and Security Concept}
\input{\TEXTprod 02_product_text17}
\section{Functional Block Overview}
\input{\TEXTprod 02_product_history_tab04}
\subsection{RISC-V Core}
\input{\TEXTprod 02_product_text18}
\subsection{Instruction Memory}
\input{\TEXTprod 02_product_text19}
\subsection{Data Memory}
\input{\TEXTprod 02_product_text20}
\subsection{Wishbone Peripherals}
\subsubsection{GPIO}
\input{\TEXTprod 02_product_text21}
\subsection{Debug}
\input{\TEXTprod 02_product_text22}

%%%%%%%%%%%%%%%%%%%%%%%%%%%%%%%%%
% C: Architecture Concepts
%%%%%%%%%%%%%%%%%%%%%%%%%%%%%%%%%
\chapter{Architecture Concepts}
\input{\TEXTarch 03_arch_text01}
\input{\TEXTarch 03_arch_bild01}
\input{\TEXTarch 03_arch_text02}
\section{Technology}
\input{\TEXTarch 03_arch_history_tab01}
\section{CPU Concept}
\input{\TEXTarch 03_arch_history_tab01a}
\input{\TEXTarch 03_arch_text03}
\input{\TEXTarch 03_arch_tab01}
\input{\TEXTarch 03_arch_tab02}
\input{\TEXTarch 03_arch_tab03}
\section{Bus Concept}
\input{\TEXTarch 03_arch_history_tab02}
\input{\TEXTarch 03_arch_bus_concept}
\section{Interrupt Concept}
\input{\TEXTarch 03_arch_history_tab03}
\section{Clock Concept}
\input{\TEXTarch 03_arch_history_tab04}
\section{System Control Concept}
\input{\TEXTarch 03_arch_history_tab05}
\section{Power Concept}
\input{\TEXTarch 03_arch_history_tab06}
\section{Memory Concept}
\input{\TEXTarch 03_arch_history_tab07}
\input{\TEXTarch 03_arch_memory_concept}
\section{System Protection and Security Concept}
\input{\TEXTarch 03_arch_history_tab08}
\section{Debug Concept}
\input{\TEXTarch 03_arch_history_tab09}
\section{Register Concept}
\input{\TEXTarch 03_arch_history_tab10}

%%%%%%%%%%%%%%%%%%%%%%%%%%%%%%%%%
% C: CPU Core
%%%%%%%%%%%%%%%%%%%%%%%%%%%%%%%%%
\chapter{RWU-RV64I Core}
\section{ALU}
\label{sec:alufunc01}
\subsubsection{History}
\input{\TEXTalu 04_alu_author_tab01}
\input{\TEXTalu 04_alu_history_tab01}
\subsubsection{Functional Configuration of ALU0}
\input{\TEXTalu 04_alu_text01}
\input{\TEXTalu 04_alu_tab01}
\subsubsection{HW Interfaces}
\input{\TEXTalu 04_alu_text02}
\input{\TEXTalu 04_alu_tab02}
\subsubsection{Bus Interface}
\input{\TEXTalu 04_alu_text03}
\subsubsection{Registers}
\input{\TEXTalu 04_alu_text04}
\subsubsection{Related Sections of this Spec}
\input{\TEXTalu 04_alu_text05}
\subsection{History of ALU}
\input{\TEXTalu 04_alu_history_tab02}
\subsection{Functional Overview}
\label{subsec:alufunc01}
\input{\TEXTalu 04_alu_text06}
\subsection{Structural Overview}
\subsubsection{I/O of the ALU}
\input{\TEXTalu 04_alu_tab03}
\subsubsection{Internals of the ALU}
%\input{\TEXTa 04_alu_tab04}
\input{\TEXTalu 04_alu_list04}
\subsection{Service Requests of the ALU}
\input{\TEXTalu 04_alu_text07}
\subsection{The ALU Peripheral Kernel}
\input{\TEXTalu 04_alu_text08}
\input{\TEXTalu 04_alu_tab05}
\input{\TEXTalu 04_alu_tab06}
\subsection{Registers}
\input{\TEXTalu 04_alu_text09}

\section{Instruction Decoder}
\label{sec:decfunc01}
\subsubsection{History}
\input{\TEXTdec 04_dec_author_tab01}
\input{\TEXTdec 04_dec_history_tab01}
\subsubsection{Functional Configuration of InstructionDecode0}
\input{\TEXTdec 04_dec_text01}
\input{\TEXTdec 04_dec_tab01}
\subsubsection{HW Interfaces}
\input{\TEXTdec 04_dec_text02}
\input{\TEXTdec 04_dec_tab02}
\subsubsection{Bus Interface}
\input{\TEXTdec 04_dec_text03}
\subsubsection{Registers}
\input{\TEXTdec 04_dec_text04}
\subsubsection{Related Sections of this Spec}
\input{\TEXTdec 04_dec_text05}
\subsection{History of Instruction Decoder}
\input{\TEXTdec 04_dec_history_tab02}
\subsection{Functional Overview}
\label{subsec:decfunc01}
\input{\TEXTdec 04_dec_text06}
\subsection{Structural Overview}
\subsubsection{I/O of the Instruction Decoder}
\input{\TEXTdec 04_dec_tab03}
\subsubsection{Internals of the Instruction Decoder}
\input{\TEXTdec 04_dec_list04}
\subsection{Service Requests of the Instruction Decoder}
\input{\TEXTdec 04_dec_text07}
\subsection{The InstrDec Peripheral Kernel}
\input{\TEXTdec 04_dec_text08}
\input{\TEXTdec 04_dec_tab05}
\subsection{Registers}
\input{\TEXTalu 04_alu_text09}

\section{Memory System}
\label{sec:memfunc01}
\input{\TEXTmems 04_memsyst_all}

%%%%%%%%%%%%%%%%%%%%%%%%%%%%%%%%%
% C: Peripherals
%%%%%%%%%%%%%%%%%%%%%%%%%%%%%%%%%
\chapter{Peripherals}
% GPIO
\section{GPIO} % see INFINEON SGOLD 2 Spec USIF p.971
\label{sec:gpiofunc01}
\subsection{System Integration of GPIO0}
\subsubsection{History}
\input{\TEXTgpio 05_gpio_author_tab01}
\input{\TEXTgpio 05_gpio_history_tab01}
\subsubsection{Functional Configuration of GPIO0}
\input{\TEXTgpio 05_gpio_text01}
\input{\TEXTgpio 05_gpio_tab01}
\subsubsection{HW Interfaces}
\input{\TEXTgpio 05_gpio_text02}
\input{\TEXTgpio 05_gpio_tab02}
\subsubsection{Bus Interface}
\input{\TEXTgpio 05_gpio_text03}
\subsubsection{Registers}
\input{\TEXTgpio 05_gpio_tab03}
\subsubsection{Related Sections of this Spec}
\input{\TEXTgpio 05_gpio_text04}
\subsection{History of GPIO}
\input{\TEXTgpio 05_gpio_history_tab02}
\subsection{Functional Overview}
\label{subsec:gpiofunc01}
\input{\TEXTgpio 05_gpio_text05}
\subsection{Structural Overview}
\input{\TEXTgpio 05_gpio_text06}
\input{\TEXTgpio 05_gpio_bild01}
\input{\TEXTgpio 05_gpio_text07}
\subsubsection{I/O of the GPIO Kernel}
\input{\TEXTgpio 05_gpio_tab04}
\subsubsection{I/O of the GPIO BPI}
\input{\TEXTgpio 05_gpio_tab05}
\subsubsection{Internals of the GPIO BPI}
\input{\TEXTgpio 05_gpio_tab07}
\subsubsection{I/O of the GPIO Top}
\input{\TEXTgpio 05_gpio_tab06}
\subsection{Service Requests of the GPIO}
\input{\TEXTgpio 05_gpio_text08}
\subsection{The GPIO Peripheral Kernel}
\label{subsec:gpiofunc02}
\input{\TEXTgpio 05_gpio_text09}
\input{\TEXTgpio 05_gpio_bild02}
\subsection{Registers}
\input{\TEXTgpio 05_gpio_text10}

% I-Mem
\section{Instruction Memory (I-Mem)}
\label{sec:imemfunc01}
\subsection{System Integration of \asimem}
\subsubsection{History}
\input{\TEXTimem 05_imem_author_tab01}
\input{\TEXTimem 05_imem_history_tab01}
\subsubsection{Functional Configuration of \asimem}
\input{\TEXTimem 05_imem_text01}
\input{\TEXTimem 05_imem_tab01}
\subsubsection{HW Interfaces}
\input{\TEXTimem 05_imem_text02}
\input{\TEXTimem 05_imem_tab02}
\subsubsection{Bus Interface}
\input{\TEXTimem 05_imem_text03}
\subsubsection{Registers}
\input{\TEXTimem 05_imem_tab03}
\subsubsection{Related Sections of this Spec}
\input{\TEXTimem 05_imem_text04}
\subsection{History of \asimemg}
\input{\TEXTimem 05_imem_history_tab02}
\subsection{Functional Overview}
\label{subsec:imemfunc01}
\input{\TEXTimem 05_imem_text05}
\subsection{Structural Overview}
\input{\TEXTimem 05_imem_text06}
%\input{\TEXTimem 05_imem_bild01}
%\input{\TEXTimem 05_imem_text07}
\subsubsection{I/O of the I-Mem Kernel}
\input{\TEXTimem 05_imem_tab04}
%\subsubsection{I/O of the I-Mem BPI}
%\input{\TEXTimem 05_imem_tab05}
%\subsubsection{Internals of the I-Mem BPI}
%\input{\TEXTimem 05_imem_tab07}
%\subsubsection{I/O of the I-Mem Top}
%\input{\TEXTimem 05_imem_tab06}
\subsection{Service Requests of the I-Mem}
\input{\TEXTimem 05_imem_text08}
\subsection{The I-Mem Peripheral Kernel}
\label{subsec:imemfunc02}
\input{\TEXTimem 05_imem_text09}
%\input{\TEXTimem 05_imem_bild02}
\subsection{Registers}
\input{\TEXTimem 05_imem_text10}

% QSPI
%\section{QSPI}
%\label{sec:qspifunc01}
%\subsection{System Integration of QSPI0}
%\subsubsection{History}
%\input{\TEXTqspi 05_qspi_author_tab01}
%\input{\TEXTqspi 05_qspi_history_tab01}
%\subsubsection{Functional Configuration of QSPI0}
%\input{\TEXTqspi 05_qspi_text01}
%\input{\TEXTqspi 05_qspi_tab01}
%\subsubsection{HW Interfaces}
%\input{\TEXTqspi 05_qspi_text02}
%\input{\TEXTqspi 05_qspi_tab02}
%\subsubsection{Bus Interface}
%\input{\TEXTqspi 05_qspi_text03}
%\subsubsection{Registers}
%\input{\TEXTqspi 05_qspi_tab03}
%\subsubsection{Related Sections of this Spec}
%\input{\TEXTqspi 05_qspi_text04}
%\subsection{History of QSPI}
%\input{\TEXTqspi 05_qspi_history_tab02}
%\subsection{Functional Overview}
%\label{subsec:qspifunc01}
%\input{\TEXTqspi 05_qspi_text05}
%\subsection{Structural Overview}
%\input{\TEXTqspi 05_qspi_text06}
%\input{\TEXTqspi 05_qspi_bild01}
%\input{\TEXTqspi 05_qspi_text07}
%\subsubsection{I/O of the QSPI Kernel}
%\input{\TEXTqspi 05_qspi_tab04}
%\subsubsection{I/O of the QSPI BPI}
%\input{\TEXTqspi 05_qspi_tab05}
%\subsubsection{Internals of the QSPI BPI}
%\input{\TEXTqspi 05_qspi_tab07}
%\subsubsection{I/O of the QSPI Top}
%\input{\TEXTqspi 05_qspi_tab06}
%\subsection{Service Requests of the QSPI}
%\input{\TEXTqspi 05_qspi_text08}
%\subsection{The QSPI Peripheral Kernel}
%\label{subsec:qspifunc02}
%\input{\TEXTqspi 05_qspi_text09}
%\input{\TEXTqspi 05_qspi_bild02}
%\subsection{Registers}
%\input{\TEXTqspi 05_qspi_text10}

%AS
\section{QSPI}
\input{\TEXTqspi 05_qspi_all}

%ChatGPT
%\input{\TEXTqspi 05_qspi_all2}

% Claude
%\section{QSPI Claude}
%\input{\TEXTqspi as_qspi_spec}

% UART
\section{UART}
\input{\TEXTuart 05_uart_all}



%\input{\TEXTqspi 05_qspi_chatgpt}

%%%%%%%%%%%%%%%%%%%%%%%%%%%%%%%%%
% C: Test & Debug
%%%%%%%%%%%%%%%%%%%%%%%%%%%%%%%%%
\chapter{Test and Debug}

\section{Simulation Test Suite}
\label{sec:simtest01}

The \texttt{RV\_NoPipeline} and \texttt{RV\_NoPipelineCache} projects include a
regression test suite driven by CMake.
Each test assembles a dedicated firmware program, runs it under xsim, and checks
the result via GPIO output (0x55\,=\,pass, 0xFF\,=\,fail) or a waveform monitor.

\subsection{UART0 Simulation Tests}
\label{subsec:uartsimtest}

Two system-level tests verify UART0 end-to-end.
Firmware uses \texttt{CLKDIV\,=\,4}, giving a bit period of
$16 \times 4 \times 800\,\text{ns} = 51\,200\,\text{ns}$ at the 1.25\,MHz core clock.

\paragraph{uart\_tx}
Firmware writes twelve ASCII bytes (\textit{``Hello Rocker''}) to the UART TX FIFO
one by one, polling STATUS[0] (tx\_busy) before each write.
The testbench monitors \texttt{uart0\_tx\_o} and decodes each 8N1 frame,
comparing against the expected string.
On success GPIO\,=\,0x55 is written.

\paragraph{uart\_rx}
The testbench injects a pre-built 8N1 serial stream
(\textit{``We will rock you.''}, 17\,bytes) onto \texttt{uart0\_rx\_i}.
Firmware polls FIFOSTAT[21] (rx\_empty) until cleared, then reads the DATA
register; each received byte is compared against the expected ASCII value.
On full match GPIO\,=\,0x55; on any mismatch GPIO\,=\,0xFF.

\subsubsection{RX FIFO Pop Bug (Fixed 2026-06-05)}
\label{subsubsec:rxfifobug}

During development, UART RX reads returned 0x00 for every received byte.
Root cause: the non-pipelined CPU (\texttt{as\_cpux}) uses a two-cycle load
sequence --- EXEC\_ST (address and read-enable asserted) followed by EXECLD\_ST
(data captured by the register file).
The Wishbone master (\texttt{asMBpi}) holds STB\,=\,CYC\,=\,1 throughout both
cycles, so the UART slave's \texttt{rd\_data\_s} signal was high during EXEC\_ST,
popping the show-ahead RX FIFO at the EXEC$\to$EXECLD clock edge.
EXECLD\_ST then saw an empty FIFO and captured 0x00.

Fix: a one-cycle delay register \texttt{rd\_data\_r} was added in
\texttt{ip\_uart\_v2/src/as\_uart\_top.sv} and used as the FIFO \texttt{rd\_en\_i}.
The pop now fires at EXECLD$\to$FETCH0, leaving the FIFO head valid throughout
EXECLD\_ST.
The same pattern must be applied to any future read-once (clear-on-read) register
on this Wishbone bus.

%%%%%%%%%%%%%%%%%%%%%%%%%%%%%%%%%
% C: Electrical Characteristics
%%%%%%%%%%%%%%%%%%%%%%%%%%%%%%%%%
\chapter{Electrical Characteristics}
Not yet

%%%%%%%%%%%%%%%%%%%%%%%%%%%%%%%%%
% C: Memory Map and Registers
%%%%%%%%%%%%%%%%%%%%%%%%%%%%%%%%%
\chapter{Memory Map and Registers}
\label{ch:memmap}

\section{System Memory Map}
\label{sec:sysmap}

Table~\ref{tab:sysmap} shows the complete address map of the \aschip{} system.
The address space is 64-bit.
Bit~32 of the address separates main memory (bit~32\,=\,0) from the peripheral
bus (bit~32\,=\,1); the peripheral bridge in \textbf{RV\_NoPipelineCache} uses
this bit to route accesses.
All peripheral registers are 64-bit wide and 8-byte aligned.
The column \textit{AW} gives the address-width parameter passed to the
peripheral's Wishbone slave interface (lower \textit{AW} bits of the full
address; the base address is always aligned to $2^{AW}$ bytes so the lower
\textit{AW} bits equal the register offset directly).

\begin{table}[H]
\caption{System Memory Map}
\label{tab:sysmap}
\centering
{\scriptsize
\begin{tabularx}{\textwidth}{|l |l |l |l |X|}
  \hline
  Base Address & Size & Peripheral & AW & Notes \\
  \hline
  \hline
  \multicolumn{5}{|l|}{\textit{Main memory (bit 32 = 0)}} \\
  \hline
  \texttt{0x0000\_0000\_0000\_0000} & 64\;kB & D-Mem / Scratchpad & --- &
    \textit{RV\_NoPipeline}: flat SRAM.
    \textit{RV\_NoPipelineCache}: scratchpad. \\
  \hline
  \multicolumn{5}{|l|}{\textit{Peripheral bus (bit 32\,=\,1)}} \\
  \hline
  \texttt{0x0000\_0001\_0000\_0000} & 512\;B & \textbf{GPIO0}  & 7 & cs[1]. Both variants. \\
  \hline
  \texttt{0x0000\_0001\_0000\_0200} & 512\;B & \textbf{QSPI0}  & 8 & cs[2]. \textit{RV\_NoPipelineCache} only. \\
  \hline
  \texttt{0x0000\_0001\_0000\_0400} & 512\;B & \textbf{CGU}    & 4 & cs[3]. Both variants. \\
  \hline
  \texttt{0x0000\_0001\_0000\_0600} & 512\;B & \textbf{UART0}  & 8 & cs[4]. Both variants. \\
  \hline
\end{tabularx}}
\end{table}

Table~\ref{tab:uartoffsets} repeats the UART0 register offsets (from the UART0
base address) for quick reference; full bit-field descriptions are in
Section~\ref{sec:uartfunc01}.

\begin{table}[H]
\caption{UART0 Register Offsets (from base \texttt{0x0000\_0001\_0000\_0600})}
\label{tab:uartoffsets}
\centering
{\footnotesize
\begin{tabularx}{\textwidth}{|l |l |l |X|}
  \hline
  Offset & Symbol & Mode & Description \\
  \hline
  \hline
  \texttt{0x00} & \asuartIDRegExt   & ro  & Peripheral ID (\texttt{0x20}) \\
  \hline
  \texttt{0x08} & \asuartLcrRegExt  & rw  & Line Control (data bits, parity, stop bits, break) \\
  \hline
  \texttt{0x10} & \asuartClkdivRegExt & rw & Baud-rate divisor (16-bit) \\
  \hline
  \texttt{0x18} & \asuartCtrlRegExt & rw  & Control (loopback, tx\_flush, rx\_flush) \\
  \hline
  \texttt{0x20} & \asuartStatusRegExt & ro & TX/RX busy; sticky error flags \\
  \hline
  \texttt{0x28} & \asuartDataRegExt & rw  & Write: TX FIFO push; Read: RX FIFO pop \\
  \hline
  \texttt{0x30} & \asuartFifostatRegExt & ro & TX/RX FIFO levels, full, half, empty \\
  \hline
  \texttt{0x38} & \asuartRxthresRegExt & rw & RX interrupt threshold (default\,=\,8) \\
  \hline
  \texttt{0x40} & \asuartRISRegExt  & ro  & Raw Interrupt Status \\
  \hline
  \texttt{0x48} & \asuartIMSCRegExt & rw  & Interrupt Mask Control \\
  \hline
  \texttt{0x50} & \asuartMISRegExt  & ro  & Masked Interrupt Status \\
  \hline
  \texttt{0x58} & \asuartICRRegExt  & wo  & Interrupt Clear (write-1-to-clear) \\
  \hline
\end{tabularx}}
\end{table}

\section{Interrupt Map}
\label{sec:irqmap}

Table~\ref{tab:irqmap} shows the assignment of external interrupt lines to the
\aschip{} CPU.

\begin{table}[H]
\caption{External Interrupt Map}
\label{tab:irqmap}
\centering
{\footnotesize
\begin{tabularx}{\textwidth}{|l |l |l |X|}
  \hline
  IRQ Bit & Source & Signal & Notes \\
  \hline
  \hline
  7 & GPIO0 & \asgpioIRQ & Input change on any GPIO pin (masked via IMSC). \\
  \hline
  6 & QSPI0 & \asqspiIRQ & DONE / RXHALF / TXHALF / ERR / TIMEOUT
    (\textit{RV\_NoPipelineCache} only). \\
  \hline
  5 & UART0 & \asuartIRQ & RX-ready, TX-ready, timeout, frame/parity/overrun/break. \\
  \hline
  4\,--\,0 & reserved & --- & Tied to 0. \\
  \hline
\end{tabularx}}
\end{table}


